\documentclass{article}
\usepackage[T1]{fontenc}
\usepackage{lmodern}
\usepackage{graphicx}
\usepackage{amsmath,amssymb}
\usepackage[a4paper,margin=2cm]{geometry}
\usepackage[absolute,overlay]{textpos}
\setlength{\TPHorizModule}{1mm}
\setlength{\TPVertModule}{1mm}
\usepackage[indonesian]{babel}
\usepackage{hyperref}
\setlength{\emergencystretch}{2em}
% unit: Halaman 3

\begin{document}

% === Halaman 3 (python/native) ===

Tinjauan Umum GOST (Magma)

• GOST = Gosudarstvenny Standard, artinya standard pemerintah, disebutkan di 

dalam standard GOST 28147-89 (RFC 5830).

• Merupakan algoritma enkripsi cipher blok standard dari negara Uni Soviet dulu 

(USSR).

• Pada mulanya cipher blok ini tidak mempunyai nama, tetapi di dalam standard 

yang baru, GOST R 34.12-2015 (RFC 7801, RFC 8891), cipher blok ini dinamakan 

Magma.

• GOST dikembangkan pada tahun 1970 namun baru dirilis kepada publik pada 

tahun 1990 setelah Uni Soviet bubar. 

• Dibuat oleh Soviet sebagai alternatif terhadap algoritma enkripsi standard 

Amerika Serikat, DES. 

• GOST secara struktural mirip dengan DES  (menggunakan jaringan Feistel)

3

\end{document}
